%! AP Statistics — Unit 5, Lesson 5.2 — Guided Notes (Answer Key)
\documentclass[10pt]{article}
\usepackage{apstats-article}
\usepackage{apstats-key}   % pulls in -boxes; do NOT also load -boxes
\newcommand{\TallMath}[1]{$\displaystyle #1\rule[-1.4em]{0pt}{3.2em}$}

\begin{document}

\pageheader{Unit 5, Lesson 5.2}{Guided Notes --- Answer Key}
\namedateperiod

\vspace{0.15in}

\begin{objectivebox}
By the end of this lesson, I will be able to:
\begin{enumerate}
  \item \ansline{find normal probabilities for an interval using z-scores or technology.}
  \item \ansline{find the boundary value for a given area (inverse normal).}
  \item \ansline{decide whether a normal model is appropriate for a distribution.}
\end{enumerate}
\end{objectivebox}

\vspace{0.1in}

\begin{vocabbox}
\noindent\textbf{\textcolor{navy}{Continuous random variable:}}\\[1pt]\ansline{A variable that can take any value in an interval;}\\[1pt]\ansline{probabilities are areas under a density curve.}\\[3pt]
\noindent\textbf{\textcolor{navy}{Normal distribution:}}\\[1pt]\ansline{A symmetric, unimodal, bell-shaped distribution}\\[1pt]\ansline{described by its mean $\mu$ and standard deviation $\sigma$.}\\[3pt]
\noindent\textbf{\textcolor{navy}{z-score:}}\\[1pt]\ansline{$z = (x-\mu)/\sigma$: the number of standard deviations}\\[1pt]\ansline{that $x$ lies above or below the mean.}\\[3pt]
\noindent\textbf{\textcolor{navy}{Standard normal distribution:}}\\[1pt]\ansline{The normal distribution with mean 0 and standard deviation 1:}\\[1pt]\ansline{$Z \sim N(0, 1)$.}\\[3pt]
\end{vocabbox}

\vspace{0.1in}

\begin{hookbox}
Adult male heights $\sim N(69.2, 3.6)$. Estimate the fraction of adult males taller
than 6 feet (72 in) from the histogram (shown on board), then we will calculate exactly.

My estimate: \ans{(varies)} \qquad Calculated answer: \ans{$\approx 0.22$ (22\%)}
\end{hookbox}

\vspace{0.1in}

\begin{notesbox}{Normal Distributions: Key Facts}
Notation: $X \sim N(\mu, \sigma)$ means $X$ is normally distributed with mean
\ans{$\mu$} and standard deviation \ans{$\sigma$}.

The \textbf{area} under the curve over an interval equals \ans{the probability for that interval} (VAR-6.A.3).

\textbf{z-score formula:} $z = $ \ans{$(x - \mu)/\sigma$} \hfill
means $z$ is the number of \ans{standard deviations} that $x$ is from the mean.

\textbf{68--95--99.7 Rule:}
\begin{itemize}
  \item $P(\mu - \sigma < X < \mu + \sigma) \approx $ \ans{0.68}
  \item $P(\mu - 2\sigma < X < \mu + 2\sigma) \approx $ \ans{0.95}
  \item $P(\mu - 3\sigma < X < \mu + 3\sigma) \approx $ \ans{0.997}
\end{itemize}
\end{notesbox}

\vspace{0.1in}

\begin{notesbox}{Calculating Normal Probabilities}
\textbf{Forward probability} (area given a boundary):
\begin{enumerate}
  \item Sketch and shade the curve. Label \ans{$\mu$} and \ans{the boundary $x$}.
  \item Compute the z-score: $z = \frac{x - \mu}{\sigma} = $ \ans{standardized boundary}
  \item Use table or \texttt{normalcdf(lower, upper, $\mu$, $\sigma$)} $=$ \ans{the probability}
  \item Interpret in context: ``The probability that \ldots'' \ansline{(context sentence)}
\end{enumerate}

\vspace{0.06in}
\textbf{Example:} $X \sim N(69.2, 3.6)$. Find $P(X > 72)$.\\[4pt]
$z = \frac{72 - 69.2}{3.6} = $ \ans{$0.78$} \qquad
$P(X > 72) = 1 - P(Z < {\color{keyred}\mathbf{0.78}}) = 1 - $ \ans{$0.7823$} $\approx$ \ans{$0.2177$}
\end{notesbox}

\vspace{0.1in}

\begin{notesbox}{Finding a Boundary Value (Inverse Normal)}
\textbf{Inverse normal} (boundary given an area):
\begin{enumerate}
  \item Sketch and shade the curve. Label the given \ans{area/percentile}.
  \item Find the z-score: $z^* = $ \ans{invNorm(area)} (from table or \texttt{invNorm}).
  \item Solve for $x$: $x = \mu + z^*\sigma = $ \ans{$\mu + z^*\sigma$ with the numbers}
\end{enumerate}

\vspace{0.06in}
\textbf{Example:} $X \sim N(69.2, 3.6)$. Find the 90th percentile.\\[4pt]
$z^* = $ \ans{$1.28$} \qquad
$x = 69.2 + ({\color{keyred}\mathbf{1.28}})(3.6) = $ \ans{$73.8$} inches
\end{notesbox}

\vspace{0.1in}

\begin{notesbox}{When Is the Normal Model Appropriate? (VAR-6.C)}
The normal model is appropriate when the distribution is \ans{approximately symmetric}
and \ans{unimodal and bell-shaped}.

Circle the correct answer for each:
\begin{itemize}
  \item Heights of adult males \hfill (\ans{appropriate} / not appropriate)
  \item Number of heads in 5 coin flips \hfill (appropriate / \ans{not appropriate})
  \item Average of 100 random observations (by CLT) \hfill (\ans{appropriate} / not appropriate)
\end{itemize}
\end{notesbox}

\vspace{0.1in}

\begin{practicebox}
$X \sim N(500, 100)$ (SAT Math scores).
\begin{enumerate}
  \item Find $P(X > 650)$. Show all steps including a sketch.

        \par\vspace{-\parskip}\nointerlineskip\noindent\begin{minipage}[t][1.0in][t]{\linewidth}\vspace*{2pt}\ans{Sketch: shade the right tail above 650.}\\
        \ans{$z = (650-500)/100 = 1.50$}\\
        \ans{$P(X>650) = 1 - 0.9332 = 0.0668$ --- about 6.7\% score above 650.}\end{minipage}

  \item Find the SAT Math score at the 25th percentile.

        \par\vspace{-\parskip}\nointerlineskip\noindent\begin{minipage}[t][0.8in][t]{\linewidth}\vspace*{2pt}\ans{$z^* = \text{invNorm}(0.25) \approx -0.674$}\\
        \ans{$x = 500 + (-0.674)(100) \approx 433$}\end{minipage}
\end{enumerate}
\end{practicebox}

\end{document}
